% !TeX program = lualatex
% !TeX encoding = utf8
% !TeX spellcheck = uk_UA
% !BIB program = biber

%% --------------------------------------------------------
\section{Функціонал енергії та натуральні орбіталі}
%% --------------------------------------------------------

%% --------------------------------------------------------
\subsection{Функціонал енергії та матриці густини в пост-HF методах}
%% --------------------------------------------------------

У методі Хартрі–Фока хвильова функція системи описується одним детермінантом
Слейтера $|\Phi_0\rangle$, який будується з спін-орбіталей $\{\psi_i\}$.
Енергія визначається як середнє значення електронного гамільтоніана,і в координатному представлені має вигляд \eqref{eq:EHForbilal_repr}.

У пост-HF методах хвильова функція вже є суперпозицією кількох
детермінантів. Наприклад в методі типу \texttt{CI}:
\[
	|\Psi\rangle = \sum_I c_I\,|\Phi_I\rangle,
\]
де $c_I$ --- коефіцієнти конфігурацій, а $|\Phi_I\rangle$ --- окремі детермінанти.
Енергія, що визначається як квантово-механічне середнє, має вигляд:
\[
	E =
	\frac{\langle \Psi | \hat{H}_\text{el} | \Psi \rangle}
	{\langle \Psi | \Psi \rangle}
	= \frac{\sum_{I,J} c_I c_J\, H_{IJ}}
	{\sum_{I,J} c_I c_J\, S_{IJ}},
\]
де $H_{IJ} = \langle \Phi_I | \hat{H}_\text{el} | \Phi_J \rangle$,
а $S_{IJ} = \langle \Phi_I | \Phi_J \rangle$
(для ортонормованих детермінантів $S_{IJ}=\delta_{IJ}$).

\paragraph{Одно- і двоелектронні густини.}
Щоб зручно виразити енергію для багатодетермінантних станів,
запроваджують \textbf{одночастинкову} та \textbf{двочастинкову матриці густини}.
%:
%\begin{align}
%\gamma_{pq} &= \langle \Psi | \hat{E}_{pq} | \Psi \rangle,
%&
%\Gamma_{pq,rs} &= \langle \Psi | \hat{E}_{pq,rs} | \Psi \rangle,
%\end{align}
%які узагальнюють поняття матриці густини \texttt{HF}.
У координатному просторі це відповідає усередненню багаточастинкової густини
за координатами всіх електронів, крім одного або двох:
\begin{align}
	\gamma(\mathbf{r}_1,\mathbf{r}_1')                            & = N
	\int \Psi^*(\mathbf{r}_1',\mathbf{r}_2,\ldots,\mathbf{r}_N)\,
	\Psi(\mathbf{r}_1,\mathbf{r}_2,\ldots,\mathbf{r}_N)\, d\mathbf{r}_2\ldots d\mathbf{r}_N, \\[4pt]
	\Gamma(\mathbf{r}_1,\mathbf{r}_2;\mathbf{r}_1',\mathbf{r}_2') & =
	N(N-1)
	\int \Psi^*(\mathbf{r}_1',\mathbf{r}_2',\ldots,\mathbf{r}_N)\,
	\Psi(\mathbf{r}_1,\mathbf{r}_2,\ldots,\mathbf{r}_N)\, d\mathbf{r}_3\ldots d\mathbf{r}_N.
\end{align}

\paragraph{Енергетичний функціонал в пост-HF методах.}
Енергію багатодетермінантного стану можна подати в орбітальному представлені у компактній формі:
\begin{equation}\label{eq:E-postHF}
	E[\gamma,\Gamma] =
	\sum_{pq} h_{pq}\, \gamma_{pq}
	+ \tfrac{1}{2}\sum_{pqrs} (pq|rs)\, \Gamma_{pq,rs} + E_0,
\end{equation}
де $(pq|rs)$ --- двоелектронні інтеграли кулонівського типу, $E_0$ --- внесок від міжядерного відштовхування та
взаємодії із зовнішніми полями. Тобто, у пост-HF методах, на відміну від \texttt{HF}, функціонал енергії є \emph{функцією двох густин} --- одно- і двочастинкової густин.

Як і у випадку метода HF, діагональні елементи $\gamma_{kk}$ дорівнюють заселеності
спін-орбіталей, а їхня сума --- повному числу електронів у молекулі (слд врахувати, що матриця перекриття детермінантів $\mathbf{S}$ --- одинична):
\begin{equation*}
	\mathrm{Tr}(\gamma) = N,\, \text{або}\, \sum\limits_{kk} \gamma_{kk} = N.
\end{equation*}

Аналогічно може бути записана в орбітальному представленні і матриця густини 2-го порядку $\Gamma_{pq,rs}$. Сума її діагональних елементів дорівнює числу пар електронів:

\begin{equation*}
	\sum_{p<q} \Gamma_{pq,pq} = \frac{N(N-1)}{2}.
\end{equation*}

У~координатному представленні елемент
$\Gamma(\mathbf{r}_1\sigma_1, \mathbf{r}_2\sigma_2; \mathbf{r}_1\sigma_1, \mathbf{r}_2\sigma_2)$
має чітке фізичне тлумачення:
він пропорційний ймовірності знайти
в~елементарному об’ємі~$d\mathbf{r}_1$ поблизу точки~$\mathbf{r}_1$
електрон зі спіном~$\sigma_1$, а~в~елементарному об’ємі~$d\mathbf{r}_2$
поблизу точки~$\mathbf{r}_2$ --- інший електрон зі спіном~$\sigma_2$,
тоді як координати та~спіни решти електронів залишаються довільними.
Інтегрування цієї функції по обох координатах приводить саме до рівності вище.

Отже, матриця густини другого (та вищих порядків) описують
\textbf{корельований рух електронів} у~молекулі.
Усі ефекти електронної кореляції зосереджені в~цих багаточастинкових матрицях густини,
тоді як матриця першого порядку~$\mathbf{\gamma}$ відображає лише середній розподіл електронів.

\begin{commentbox}
	У разі однодетермінантного наближення (метод Хартрі-Фока) матриця густини 2-го порядку виражається
	через матрицю густини 1-го порядку. Це означає, що вона не містить ніякої нової інформації в порівнянні з
	матрицею густини 1-го порядку. Це і логічно, бо в методі Хартрі-Фока ефекти, пов'язані з електронною кореляцією, не
	враховуються.
\end{commentbox}

Проведемо розрахунок енергії молекули \ce{N2} методом \texttt{CI} двома способами: спочатку це зробимо прямо через API \PySCF, а також по формулі \eqref{eq:E-postHF}, використовуючи матриці густини $\gamma$ та $\Gamma$ та використовуючі високорівневі API:
\begin{minted}{python}
from pyscf import gto, scf, ci
import numpy as np

mol = gto.M(atom='N 0 0 0; N 0 0 1.1', basis='cc-pVDZ')
mf = scf.RHF(mol).run()

# Спосіб 1
mci = ci.CISD(mf).run()

# Отримуємо матриці густини
rdm1 = mci.make_rdm1()
rdm2 = mci.make_rdm2()

# Отримуємо інтеграли
h1e = mf.get_hcore()
h2e = mf._eri
nuc_energy = mol.energy_nuc()
mo_coeff = mf.mo_coeff

# Спосіб 2

# Перетворюємо все в MO базис
print("Обчислення в MO базисі")
h1e_mo = np.einsum('pi,pq,qj->ij', mo_coeff, h1e, mo_coeff)

# Перетворюємо 2e-інтеграли в MO базис (це може бути повільно для великих систем)
# Використовуємо pyscf функцію для ефективного перетворення
from pyscf import ao2mo
h2e_mo = ao2mo.full(mf._eri, mo_coeff)

# Конвертуємо в повну 4-індексну матрицю
h2e_mo = ao2mo.restore(1, h2e_mo, mo_coeff.shape[1])

# Обчислюємо енергію в MO базисі за формулою
# E = Σ h1e_ij * rdm1_ij + 0.5 * Σ h2e_ijkl * rdm2_ijkl + E_nuc
energy_mo = (np.einsum('ij,ji->', h1e_mo, rdm1) +
             0.5 * np.einsum('ijkl,ijkl->', h2e_mo, rdm2) +
             nuc_energy)

print(f"Енергія CISD: {mci.e_tot}")
print(f"Енергія з RDM (MO): {energy_mo}")
print(f"Різниця: {abs(mci.e_tot - energy_mo):.10e}")
print(f"Чи однакові? {np.allclose(mci.e_tot, energy_mo)}\n")
\end{minted}
Результати роботи коду показують, що обчислення однакові в межах розрахункової похибки.


%\paragraph{Порівняння методів.}
%\begin{center}
%\begin{tabular}{lll}
%\toprule
%Характеристика & \texttt{HF} & \texttt{CI} \\
%\midrule
%Тип хвильової функції & один детермінант & суперпозиція детермінантів \\
%Функціонал енергії & $E[\mathbf{D}]$ & $E[\gamma,\Gamma]$ \\
%Варіювання & по орбіталях & по коефіцієнтах $c_I$ \\
%Кореляція & відсутня & врахована явно \\
%Основне рівняння & $F C = S C \varepsilon$ & $H C = S C E$ \\
%\bottomrule
%\end{tabular}
%\end{center}

%% --------------------------------------------------------
\subsection{Натуральні орбіталі і заселеності}
%% --------------------------------------------------------

Отримана в~пост-Хартрі–Фок розрахунках матриця густини 1-го порядку $\gamma$, як правило, не є діагональною.
Діагоналізація цієї матриці визначає нову систему орбіталей, які називають \emph{натуральними орбіталями}:
\[
	\gamma\, \phi_i = n_i\, \phi_i,
\]
де $\phi_i$ --- натуральні орбіталі, а $n_i$ --- власні значення матриці густини $n_i$ --- це заселеності відповідних натуральних орбіталей,
що можуть набувати будь-яких значень між $0 \leqslant 1 \leqslant 2$.

П.~О.~Льовдін показав, що в~базисі натуральних орбіталей
розклад хвильової функції за детермінантами має найкращу збіжність:
кількість суттєвих  конфігурацій у~ньому є мінімальною.
Завдяки цьому натуральні орбіталі забезпечують
найкомпактніше та~найнаочніше представлення електронної структури атомів і~молекул,
а~також широко використовуються для аналізу результатів \texttt{CI}- та~\texttt{CASSCF}-розрахунків.

\paragraph{Практична цінність.}
\begin{itemize}
	\item У базисі натуральних орбіталей хвильова функція має найкомпактніший вигляд ---
	      більшість малозначущих детермінантів можна відкинути без втрати точності.
	\item Числа зайнятості $n_i$ дозволяють оцінити силу електронної кореляції:
	      чим ближчі вони до 1, тим сильніше система відхиляється від однодетермінантного опису.
	\item У багатоконфігураційних методах (\texttt{CASSCF}) натуральні орбіталі
	      допомагають ідентифікувати \emph{активний простір} --- тобто ті орбіталі,
	      які реально беруть участь у збудженнях і хімічних процесах.
\end{itemize}

У~скрипті \inputcode{h2o_casscf_natural_orbitals.py}
реалізуємо програму, яка дозволить
послідовно виконати \texttt{RHF}-, \texttt{CISD}- та~\texttt{CASSCF}-розрахунки,
отримати матрицю густини першого порядку,
побудувати натуральні орбіталі та проаналізувати їхні заселеності.
На підставі цього можна автоматично визначити
фракційно заселені орбіталі й сформувати оптимальний
активний простір для подальших \texttt{CASSCF} розрахунків. Далі ми це покажемо покроково:

\begin{enumerate}
	\item \textbf{RHF~--- початкові орбіталі.}
	      Виконується стандартний розрахунок середнього поля, що дає молекулярні орбіталі Хартрі–Фока:
	      \begin{minted}{python}
mf = scf.RHF(mol).run()
	\end{minted}

	\item \textbf{CISD~--- урахування електронної кореляції.}
	      На базі отриманих орбіталей проводиться конфігураційний розклад із збудженнями до подвійних:
	      \begin{minted}{python}
myci = ci.CISD(mf).run()
	\end{minted}
	      Цей крок враховує кореляцію електронів і формує більш фізично коректну однопартіційну матрицю густини (1-RDM).

	\item \textbf{Діагоналізація матриці густини.}
	      Власні значення 1-RDM відповідають заселеностям натуральних орбіталей~$n_i$:
	      \begin{minted}{python}
dm1 = myci.make_rdm1()
occs, natorbs = np.linalg.eigh(dm1)
	\end{minted}
	      Діагоналізація перетворює початковий орбітальний базис у базис натуральних орбіталей.

	\item \textbf{Вибір активного простору.}
	      Орбіталі з фракціональними заселеностями
	      $0.02 < n_i < 1.98$ включаються до активного простору~\texttt{CAS}:
	      \begin{minted}{python}
active = [i for i, n in enumerate(occs) if 0.02 < n < 1.98]
ncas = len(active)
nelecas = int(round(sum(occs[i] for i in active)))
	\end{minted}

	\item \textbf{CASSCF у базисі натуральних орбіталей.}
	      З отриманим активним простором проводиться багатоконфігураційний розрахунок:
	      \begin{minted}{python}
mf.mo_coeff = mf.mo_coeff @ natorbs
mc = mcscf.CASSCF(mf, ncas, nelecas)
mc.kernel()
	\end{minted}
	      Розрахунок у базисі натуральних орбіталей зазвичай забезпечує найкращу збіжність і найкомпактніше представлення хвильової функції.

	\item \textbf{Аналіз результатів.}
	      Після завершення \texttt{CASSCF} можна знову обчислити заселеності, щоб оцінити якість вибору активного простору:
	      \begin{minted}{python}
dm1_cas = mc.make_rdm1()
occs_cas, _ = np.linalg.eigh(dm1_cas)
	\end{minted}
\end{enumerate}

\begin{commentbox}
	У наведеному коді натуральні орбіталі отримані
	через явне обчислення одночастинкової матриці густини
	та її діагоналізацію:
	\inlinecode{dm1 = mc.make\_rdm1()} і \inlinecode{np.linalg.eigh(dm1)}.
	Однак у~\texttt{PySCF} це саме можна зробити коротше
	за допомогою вбудованої API-функції\\  \inlinecode{addons.make\_natural\_orbitals(mc)},
	яка автоматично виконує ці дії та повертає пари
	\texttt{(natorb, occ)}~--- натуральні орбіталі й відповідні заселеності:
	\begin{minted}{python}
from pyscf.mcscf import addons
natorb, occ = addons.make_natural_orbitals(mc)
print(occ)
\end{minted}
\end{commentbox}

%Функція \inlinecode{addons.make\_natural\_orbitals(mc)}
%діагоналізує одночастинкову матрицю густини
%та повертає пари \texttt{(natorb, occ)} ---
%коефіцієнти орбіталей і відповідні їм числа зайнятості.
%Ці орбіталі можна використовувати як оптимізований базис
%для наступних багатоконфігураційних або пост-Хартрі–Фоківських розрахунків.

%\paragraph{Примітка.}
%Старий атрибут \inlinecode{mycas.natorb = True} більше не використовується.
%Він лише обертав орбіталі в активному просторі,
%не створюючи повної натуральної бази.
%Новий інтерфейс через \inlinecode{addons.make\_natural\_orbitals()}
%дає коректну густину та справжні натуральні орбіталі для всієї системи.

%\paragraph{Типові значення чисел зайнятості.}
%
%Числа зайнятості натуральних орбіталей $n_i$ наочно показують,
%наскільки багатодетермінантною є хвильова функція.
%У табл.~\ref{tab:natorb_occup} наведено характерні значення $n_i$
%для систем із різним рівнем електронної кореляції.

%\begin{table}[h!]
%	\centering
%	\caption{Типові значення чисел зайнятості $n_i$ для різних режимів електронної кореляції.}
%	\label{tab:natorb_occup}
%	\begin{tblr}{
%		colspec={Q[c,m]Q[c,m]Q[c,m]X[l,m]},
%		row{1} = {font=\bfseries, bg=gray9},
%		hlines, vlines
%		}
%		\SetCell[r=2]{c} Рівень опису   &
%		\SetCell[r=2]{c} Тип системи    &
%		\SetCell[r=2]{c} Приклади $n_i$ &
%		\SetCell[r=2]{l} Характеристика                                                                                                                                              \\
%
%		% ----- data rows -----
%		\texttt{RHF}                    & Однодетермінантна & 2.00, 2.00, 0.00, 0.00 & Електрони повністю локалізовані у зайнятих орбіталях. Кореляція відсутня.                     \\
%		\texttt{MP2}, \texttt{CCSD}     & Слабка кореляція  & 1.98, 1.95, 0.05, 0.02 & Малий обмін густини між орбіталями. Відхилення від 2/0 --- кілька відсотків.                  \\
%		\texttt{CASSCF}                 & Помірна кореляція & 1.90, 1.75, 0.25, 0.10 & Активні орбіталі частково заселені; хвильова функція змішана.                                 \\
%		\texttt{CASSCF}, \texttt{DMRG}  & Сильна кореляція  & 1.2, 1.0, 0.8, 0.6     & Електрони значно делокалізовані між кількома орбіталями; однодетермінантний опис непридатний. \\
%	\end{tblr}
%\end{table}

Як видно, чим далі $n_i$ від 2 або 0, тим сильніше проявляється електронна кореляція.
Набір натуральних орбіталей і чисел зайнятості таким чином забезпечує
наочне уявлення про структуру хвильової функції
та допомагає оцінити, які орбіталі слід включати до активного простору.

